% translate_docs: ./manuals/car-hacking_manual.tex @ 0a5b5818922daeacebc425136aa5f5af3917ba82
\documentclass[superscriptaddress,onecolumn,floatfix]{revtex4-2}

\usepackage{graphicx}
\usepackage{epsfig}
\usepackage{color}
\usepackage{amssymb}
\usepackage{amsmath}
\usepackage{array}
\usepackage[loose]{units}
\usepackage{hyperref}
\usepackage{cleveref}
\usepackage{tasks}
\usepackage[caption=false]{subfig}
\usepackage[german]{babel}
\usepackage{letltxmacro}

\LetLtxMacro{\ORIGselectlanguage}{\selectlanguage}
\makeatletter
\DeclareRobustCommand{\selectlanguage}[1]{%
    \@ifundefined{alias@\string#1}
      {\ORIGselectlanguage{#1}}
      {\begingroup\edef\x{\endgroup
         \noexpand\ORIGselectlanguage{\@nameuse{alias@#1}}}\x}%
}
\newcommand{\definelanguagealias}[2]{%
  \@namedef{alias@#1}{#2}%
}
\makeatother

\definelanguagealias{de}{german}


\begin{document}
\selectlanguage{de}

\begin{abstract}
  Dieses Handbuch behandelt die grundlegende Verwendung und häufig gestellte Fragen zum Car-Hacking-Kit.
\end{abstract}
\title{Handbuch für Auto-Hacking-Kits}
\author{Electroniq Buttons Boutique}
\email{support@electroniqbuttons.com}
\date{\today}

\maketitle


\section{Welche Autos können das nutzen?} \label{sect:eligibility}
Das Kit ist für die meisten E-GMP-Fahrzeuge geeignet. Der ICU-H-Kabelbaum ist für die Modelle Ioniq 5s (2021–2024) und EV6s (2022–2024) vorgesehen. Der CCU-Kabelbaum ist für die Modelle Ioniq 6s und GV60s (2023–2025), EV9s (2024–2027), Ioniq 5s und EV6s (2025–2027), Ioniq 9s (2026–2027) sowie Kona und Niro EVs (2024–2026) geeignet. Andere Modelle oder Modellaktualisierungen, die ab 2022 auf den Markt kamen, sind wahrscheinlich ebenfalls kompatibel. \href{mailto:requests@electroniqbuttons.com}{Kontakt aufnehmen} Wenn Sie Interesse haben, prüfen wir die Kompatibilität. 

\section{Grundausstattung}
Beginnen Sie mit dem \href{https://github.com/tylerharvey/Ioniq5_CAN/blob/main/guides/}{Installationsanleitungen}. Sehen \href{https://meatpihq.github.io/wican-fw/config/wifi}{WiCAN-Dokumente} Für die grundlegende WiCAN-Konfiguration. Dies ist für Reverse Engineering nicht zwingend erforderlich, wird aber aus Sicherheitsgründen dringend empfohlen.

\section{SavvyCAN}
Es gibt viele Möglichkeiten, CAN-Protokolle mit WiCAN zu erfassen, aber die meisten Leute wählen \href{https://www.savvycan.com/}{SavvyCAN} Für einfache Bedienung und optimale Funktionen. Eine kurze Kurzanleitung ist verfügbar. \href{https://meatpihq.github.io/wican-fw/savvycan/usage}{Hier}Wir hoffen, Ihnen bald mehr anbieten zu können.

Ein Fork von SavvyCAN, der aktiv von Mitgliedern dieser E-GMP-Autohacking-Community entwickelt wird. \href{https://github.com/SuperSuave/SavvyLens}{SavvyLens}, verfügt über Funktionen, die speziell für unseren Reverse-Engineering-Kontext entwickelt wurden, wie z. B. das Setzen von Signal-Lesezeichen und einen benutzerfreundlicheren Signal-Viewer.

\section{Ihr Auto zurücksetzen}

Das Senden von CAN-Frames an Ihr Fahrzeug beim Reverse Engineering kann zu Fehlfunktionen des Steuergeräts führen. Wir haben Probleme mit der Klimabedienung und der Helligkeitsregelung beim Reverse Engineering dieser Steuergeräte festgestellt. Falls Ihr Fahrzeug ungewöhnlich reagiert, trennen Sie die 12-V-Stromversorgung für 30 Minuten, um einen Hard-Reset durchzuführen. Hinweis: Dabei gehen die BMS-Daten verloren. Sichern Sie diese daher unbedingt vorher.

\section{Häufig gestellte Fragen}
\begin{enumerate}
  \item[Q:] Warum werden mir keine CAN-Daten angezeigt?
  \item[A:] Dafür kann es viele Gründe geben. WiCAN lässt sich so konfigurieren, dass dies passiert (z. B. mit einem anderen Protokoll als „savvycan“), die Standardeinstellungen sollten aber funktionieren. Manche Busse verwenden CAN-FD, und mit einem CAN 2.0 WiCAN werden bei diesen keine Daten angezeigt. Wenn Ihr Auto länger als 5 Minuten ausgeschaltet und verriegelt ist, werden ebenfalls keine Daten angezeigt. Probleme mit der WLAN-Verbindung können ebenfalls die Ursache sein. Prüfen Sie daher, ob in SavvyCAN noch ein Bus verbunden ist.
  \item[Q:] Wie erhalte ich Zugriff auf CAN-FD-Busse?
  \item[A:] Wir planen, demnächst ein WiCAN mit CAN-FD anzubieten. Bis dahin ist für alle, die es eilig haben, der Lilygo T-2CAN-FD in Kombination mit einem OBD-Stecker von AliExpress eine günstige Alternative.
\end{enumerate}
% \begin{figure}[h]
%   \subfloat[]{\label{subfig:overview:harness}
%   \includegraphics[height=1.95in]{photos/whole_harness.jpg} }
%   \subfloat[]{\label{subfig:overview:shunt_cap_labeled}
%   \includegraphics[height=1.95in]{photos/labeled_connectors_MITM_harness.jpg} }
%    \caption{(a) Man-in-the-middle harness. (b) Overview of connectors showing shunt cap. \label{fig:overview}}
% \end{figure}
\appendix 
\bibliographystyle{apsrev4-1}
\bibliography{stemh_model}{}
\end{document}
